\documentclass[11pt,twoside]{article}

\usepackage[margin=1in]{geometry}
\usepackage{amsmath,amssymb,amsthm,mathtools}
\usepackage[inline,shortlabels]{enumitem}
\usepackage{xcolor}
\usepackage{microtype}

% Build identity: keep this block in future lecture notes as well.
% The timestamp is fixed at compilation start and appears on every page.
\usepackage[showseconds=false,showzone=true]{datetime2}
\DTMsetup{datesep=-,datetimesep={\space},timezonesep={\space UTC}}
\usepackage{eso-pic}
\AddToShipoutPictureBG{%
  \AtPageLowerLeft{%
    \put(\LenToUnit{0.5\paperwidth},\LenToUnit{0.5in}){%
      \makebox(0,0){\normalfont\footnotesize\color{black!60}Compiled: \DTMnow}%
    }%
  }%
}

\usepackage[colorlinks=true,linkcolor=blue!60!black,
  urlcolor=blue!60!black]{hyperref}
\usepackage[nameinlink,capitalize]{cleveref}
\crefname{exercise}{Exercise}{Exercises}
\Crefname{exercise}{Exercise}{Exercises}

\definecolor{courseblue}{HTML}{17365D}

\newcommand{\I}{\mathbb I}
\newcommand{\E}{\mathbb E}
\renewcommand{\P}{\mathbb P}
\newcommand{\R}{\mathbb R}
\makeatletter
\@for\@letter:=A,B,C,D,E,F,G,H,I,J,K,L,M,N,O,P,Q,R,S,T,U,V,W,X,Y,Z\do{%
  \expandafter\edef\csname c\@letter\endcsname{\noexpand\mathcal{\@letter}}}
\makeatother

\newtheorem{theorem}{Theorem}[section]
\newtheorem{proposition}[theorem]{Proposition}
\newtheorem{corollary}[theorem]{Corollary}
\theoremstyle{definition}
\newtheorem{definition}[theorem]{Definition}
\theoremstyle{remark}
\newtheorem{remark}[theorem]{Remark}

\renewcommand{\thepage}{3--\arabic{page}}
\renewcommand{\thesection}{3.\arabic{section}}
\renewcommand{\theequation}{3.\arabic{equation}}

\pagestyle{myheadings}
\markboth{Lecture 3: Probabilistic models, log loss, and compression}
{Lecture 3: Probabilistic models, log loss, and compression}

\usepackage{../exercises/course-exercises}

\begin{document}

\begin{center}
\fcolorbox{courseblue}{white}{%
\begin{minipage}{0.92\textwidth}
\vspace{1mm}
\textbf{CMPUT 654: Mathematical Foundations of Modern AI Systems}
\hfill Fall 2026

\vspace{3mm}
\begin{center}
{\Large\bfseries Lecture 3: Probabilistic models, log loss, and compression}

\vspace{1mm}
September 8, 2026
\end{center}

\vspace{2mm}
\textit{Lecturer: Csaba Szepesv\'ari}
\hfill
\textit{Scribe: Student name}
\vspace{1mm}
\end{minipage}}
\end{center}

% Add the scribe note above the exercises. A useful initial structure is:
% \section{Probabilistic prediction and log loss}
% \section{Maximum likelihood and cross-entropy}
% \section{KL divergence and compression}
% \section{Take-home messages}
% \section{Glossary}
% \section*{Bibliographic remarks}

\section{Exercises}

\exerciseinstructions
\setcounter{courseexercise}{0}
% Shared exercise chunk: l3-probabilistic-classification
\begin{exercisechunk}
\item \textbf{Classification from probabilistic prediction.}
\label[exercise]{ex:probabilistic-classification}
Let $(X,Y)$ have an arbitrary distribution on $\cX\times\{0,1\}$, where
$\cX$ is finite. For every $x$ with $\P(X=x)>0$, define
\[
\eta(x)=\P(Y=1\mid X=x).
\]
At the remaining inputs, $\eta(x)$ may be defined arbitrarily. For a decision
$a\in\{0,1\}$ and an $x$ with $\P(X=x)>0$, define the conditional risk
\[
r_x(a)=\P(a\ne Y\mid X=x).
\]
For a classifier $h:\cX\to\{0,1\}$, define its risk by
\[
R_{01}(h)=\P(h(X)\ne Y).
\]
A classifier is \emph{Bayes-optimal} if it minimizes $R_{01}$ over all
classifiers $h:\cX\to\{0,1\}$. The minimum value
\[
R_{01}^*=\min_{h:\cX\to\{0,1\}}R_{01}(h)
\]
is the \emph{Bayes risk}.
\begin{enumerate}[(a)]
\item Compute $r_x(0)$ and $r_x(1)$. Show that
\[
h^*(x)=\I\{\eta(x)\ge1/2\}
\]
is Bayes-optimal and that the Bayes risk is
\[
R_{01}^*=\mathbb E[\min\{\eta(X),1-\eta(X)\}].
\]
\item Prove the exact excess-risk identity
\[
R_{01}(h)-R_{01}^*
=\mathbb E\left[
|2\eta(X)-1|\I\{h(X)\ne h^*(X)\}
\right].
\]
\item Let $q:\cX\to(0,1)$ be a probabilistic
predictor, and define
\[
L_{\log}(q)
=\mathbb E[-Y\log q(X)-(1-Y)\log(1-q(X))].
\]
Let $L_{\log}^*=L_{\log}(\eta)$, with the usual conventions at
$\eta(x)\in\{0,1\}$, and let $h_q(x)=\I\{q(x)\ge1/2\}$. Prove
\[
R_{01}(h_q)-R_{01}^*
\le \sqrt{2\bigl(L_{\log}(q)-L_{\log}^*\bigr)}.
\]
\item \emph{Optional extension.} Construct a family of probabilistic
predictors whose induced
classifiers $h_q$ are Bayes-optimal but whose excess log loss tends to
infinity.
\end{enumerate}

\emph{Hint for (c):} You may use Pinsker's inequality for Bernoulli
distributions:
\[
|p-q|\le
\sqrt{\frac12\left(
p\log\frac pq+(1-p)\log\frac{1-p}{1-q}
\right)}.
\]

\paragraph{What this exercise is meant to teach.}
A probabilistic predictor and a classifier answer different questions. The
excess classification risk weights a wrong decision by twice the distance of
$\eta(x)$ from the decision boundary $1/2$. Small excess log loss controls
excess classification risk, but Bayes-optimal classification does not imply
accurate probabilities or small log loss.
\end{exercisechunk}

% Shared exercise chunk: l3-fresh-test-evaluation
\begin{exercisechunk}
\item \textbf{Evaluation on a fresh test set.}
\label[exercise]{ex:fresh-test-evaluation}
Let $\cX$ be finite. Let $D=((X_i,Y_i))_{i=1}^m$ be an i.i.d. test sample
from a distribution on $\cX\times\{0,1\}$, and let $(X,Y)$ be an independent
draw from the same distribution. For a classifier $h:\cX\to\{0,1\}$, define
\[
R(h)=\P(h(X)\ne Y),
\qquad
\widehat R_D(h)=\frac1m\sum_{i=1}^m\I\{h(X_i)\ne Y_i\}.
\]
\begin{enumerate}[(a)]
\item Show that every classifier $h:\cX\to\{0,1\}$ satisfies, for every
$\epsilon>0$,
\[
\P\bigl(|\widehat R_D(h)-R(h)|>\epsilon\bigr)
\le 2e^{-2m\epsilon^2}.
\]
\item Fix a finite set $\cH=\{h_1,\ldots,h_K\}$ of classifiers. Use the
union bound to prove
\[
\P\left(\max_{h\in\cH}|\widehat R_D(h)-R(h)|>\epsilon\right)
\le 2K e^{-2m\epsilon^2}.
\]
Give a sufficient sample size for the displayed maximum to be at most
$\epsilon$ with probability at least $1-\delta$.

\item
Let $f: (\cX\times \{0,1\})^m \to \cH$ and let
$\widehat h=f(D)$.
Show that
 \[
\P\left(
\left|\widehat R_D(\widehat h)-R(\widehat h)\right|>\epsilon
\right)
\le 2K e^{-2m\epsilon^2}.
\]

\item
The previous part uses a simultaneous guarantee over all classifiers that the
selection rule may return. Show what can happen without such a guarantee: a
classifier defined using the test sample can have
zero test error but large population risk.
In particular, consider the following:
Let $\cX=[N]$, let $X$ be uniform on $[N]$, and let $Y=1$
almost surely. Define the following function of the sample $D$:
\[
h_D(x)=\I\{x\in\{X_1,\ldots,X_m\}\}.
\]
Compute $\widehat R_D(h_D)$ and $R(h_D)$. Show that $N\ge2m$ implies
$R(h_D)\ge1/2$, and explain why this does not contradict part~(a) or part~(b).
\end{enumerate}
\emph{Hint for (a):} Apply Hoeffding's inequality to the independent random
variables $\I\{h(X_i)\ne Y_i\}$.
\end{exercisechunk}


\end{document}
